\documentclass[11pt]{article}
\usepackage[margin=1in]{geometry}
\usepackage{amsmath,amssymb,amsthm}
\usepackage{graphicx}
\usepackage[hidelinks]{hyperref}

\newtheorem{theorem}{Theorem}
\newtheorem{lemma}{Lemma}
\newtheorem{remark}{Remark}

\newcommand{\R}{\mathbb R}
\newcommand{\Id}{\mathrm{Id}}
\newcommand{\conv}{\operatorname{conv}}
\newcommand{\spanop}{\operatorname{span}}
\newcommand{\tr}{\operatorname{tr}}
\newcommand{\relint}{\operatorname{relint}}

\title{A Direct Linear-Algebra Proof of the Carlen--Cordero-Erausquin Radial-Isotropic Criterion}
\author{Run \texttt{fa\_banach\_001}, lane 16}
\date{2026-07-18}

\begin{document}
\maketitle

\section*{Status and Source}

This packet answers the proof-style question in Carlen--Cordero-Erausquin,
\emph{Subadditivity of the Entropy and its Relation to Brascamp--Lieb Type
Inequalities}, arXiv:0710.0870 \cite{CCE}.  On page 30, just before Theorem
4.3, the authors write that the theorem concerns a problem in linear algebra
and that they do not know a direct proof in a purely linear algebra context.
The theorem itself is already proved in the source via entropy and
Brascamp--Lieb extremizers.  The contribution here is a direct proof of the
same linear-algebra criterion.

\begin{center}
\includegraphics[width=\textwidth]{figures/open_problem_crop.png}
\end{center}

\section*{Source Question}

Let \(A=\{a_1,\ldots,a_m\}\) span \(\R^n\), \(n\ge 2\), and let
\(0\le c_j\le 1\).  The source asks for a direct linear-algebra proof of the
criterion for the existence of an invertible symmetric matrix \(R\) such that
\[
  \sum_{j=1}^m c_j
  \left(\frac{Ra_j}{|Ra_j|}\right)\otimes
  \left(\frac{Ra_j}{|Ra_j|}\right)
  =\Id_{\R^n}.
  \tag{1}
\]
The criterion is that \(A\) is totally reducible for \(c\), in the sense of the
source's Definition 3.8.

\section*{Definitions}

For \(J\subseteq\{1,\ldots,m\}\), write
\[
  V_J=\spanop\{a_j:j\in J\}.
\]
The Brascamp--Lieb rank polytope associated with \(A\) is
\[
  K_A=
  \left\{
    c\in[0,1]^m:
    \sum_{j=1}^m c_j=n,\quad
    \sum_{j\in J}c_j\le \dim V_J \text{ for all }J
  \right\}.
\]
Its strict part is
\[
  K_A^\circ=
  \left\{
    c\in K_A:
    \sum_{j\in J}c_j<\dim V_J
    \text{ for all nonempty proper }J
  \right\}.
\]
After discarding zero-weight indices, \(A\) is \emph{totally reducible for
\(c\)} if the remaining index set can be partitioned into nonempty sets
\(J_1,\ldots,J_r\) such that
\[
  \R^n=V_{J_1}\oplus\cdots\oplus V_{J_r},
  \qquad
  c_{J_\ell}\in K_{A_{J_\ell}}^\circ
  \quad(\ell=1,\ldots,r),
\]
where each \(A_{J_\ell}\) is considered inside \(V_{J_\ell}\).

\section*{Proof Intuition}

The equation is exactly the condition that a linear image of the vectors is in
weighted radial isotropic position.  For the strict/interior case, one can find
the required change of variables by minimizing a log-det function.  Cauchy--
Binet rewrites this function as a log-sum-exp over the bases of the vector
matroid of \(A\).  The rank inequalities defining \(K_A^\circ\) are precisely
the statement that \(c\) lies in the relative interior of the basis polytope,
which makes the log-sum-exp coercive modulo its harmless lineality.  At the
minimum, the leverage-score equations give (1).

For boundary points, every saturated rank inequality forces an actual
orthogonal splitting after any solution of (1): taking the trace against the
orthogonal projection onto the saturated span leaves no room for positive
weight vectors outside that span to have a component inside it.  Conversely,
if total reducibility gives a direct-sum decomposition into strict blocks, solve
the strict problem on each block and assemble the block solutions.  A final
polar-decomposition step makes the resulting map symmetric positive definite.

\section*{Main Result}

\begin{theorem}
Let \(A=\{a_1,\ldots,a_m\}\) span \(\R^n\), \(n\ge2\), and let
\(0\le c_j\le1\).  Ignoring indices with \(c_j=0\), assume all positive-weight
vectors are nonzero.  Then there exists an invertible symmetric matrix \(R\)
satisfying \((1)\) if and only if \(A\) is totally reducible for \(c\).
\end{theorem}

\begin{proof}
We give a proof independent of entropy and Brascamp--Lieb extremizers.

\medskip
\noindent\textbf{Step 1: the strict case.}
Let \(\mathcal B\) denote the family of \(n\)-element subsets
\(B\subseteq\{1,\ldots,m\}\) for which \(\{a_j:j\in B\}\) is a basis of
\(\R^n\).  By the vector-matroid base-polytope theorem of Edmonds
\cite{Edmonds}, the polytope
\[
  \conv\{1_B:B\in\mathcal B\}
\]
is exactly \(K_A\).  Thus \(c\in K_A^\circ\) is the same as saying that \(c\)
lies in the relative interior of this basis polytope.

Assume \(c\in K_A^\circ\).  For \(t=(t_1,\ldots,t_m)\in\R^m\), put
\[
  M(t)=\sum_{j=1}^m e^{t_j}a_j\otimes a_j
\]
and define
\[
  F(t)=\log\det M(t)-\sum_{j=1}^m c_jt_j.
\]
By Cauchy--Binet,
\[
  \det M(t)=
  \sum_{B\in\mathcal B} \det(A_B)^2
  \exp\left(\sum_{j\in B}t_j\right),
\]
where \(A_B\) is the \(n\times n\) matrix with columns \(a_j\), \(j\in B\).
Therefore
\[
  F(t)=
  \log\sum_{B\in\mathcal B}\det(A_B)^2
     e^{\langle 1_B-c,t\rangle}.
\]
Since \(0\in\relint\conv\{1_B-c:B\in\mathcal B\}\), this log-sum-exp is
coercive on the quotient by its lineality space and hence attains its minimum.
Let \(t^*\) be a minimizer.  Differentiating at \(t^*\) gives, for every \(j\),
\[
  c_j=e^{t_j^*}\,\langle M(t^*)^{-1}a_j,a_j\rangle.
  \tag{2}
\]
Set \(R=M(t^*)^{-1/2}\), the positive square root.  Then
\[
\begin{aligned}
 \sum_{j=1}^m c_j
 \left(\frac{Ra_j}{|Ra_j|}\right)\otimes
 \left(\frac{Ra_j}{|Ra_j|}\right)
 &= \sum_{j=1}^m
    e^{t_j^*}\, Ra_j\otimes Ra_j \\
 &= R\left(\sum_{j=1}^m e^{t_j^*}a_j\otimes a_j\right)R
  = \Id.
\end{aligned}
\]
Thus the desired matrix exists in the strict case.

\medskip
\noindent\textbf{Step 2: a solution forces splitting on saturated sets.}
Suppose now that an invertible \(R\) satisfying (1) exists.  Put
\[
  u_j=\frac{Ra_j}{|Ra_j|}
\]
for all \(j\) with \(c_j>0\).  Then
\[
  \sum_{j:c_j>0}c_j\,u_j\otimes u_j=\Id.
  \tag{3}
\]
Taking traces in (3) gives \(\sum_jc_j=n\).  If \(J\) is any set of positive
weight indices and \(P_J\) is the orthogonal projection onto
\(\spanop\{u_j:j\in J\}\), then
\[
  \dim\spanop\{u_j:j\in J\}
  =\tr P_J
  =\sum_i c_i |P_Ju_i|^2
  \ge \sum_{j\in J}c_j.
  \tag{4}
\]
Because \(R\) is invertible, this is exactly the rank inequality for the
original vectors \(a_j\).  Hence \(c\in K_A\).

Moreover, equality in (4) for a nonempty proper \(J\) is rigid.  Since
\(|P_Ju_j|=1\) for \(j\in J\), equality can hold only if
\[
  P_Ju_i=0
  \qquad(i\notin J,\ c_i>0).
  \tag{5}
\]
Thus every positive-weight vector outside a saturated span is orthogonal to
that span, after applying \(R\).

If there is no nonempty proper saturated set among the positive-weight indices,
then \(c\in K_A^\circ\) and \(A\) is totally reducible with one block.  Otherwise
choose an inclusion-minimal nonempty saturated set \(J\).  Minimality implies
\(c_J\in K_{U_J}^\circ\), where \(U_J=\{u_j:j\in J\}\) is considered in its
span.  By (5), the remaining positive-weight vectors lie in the orthogonal
complement of \(\spanop U_J\), and (3) splits into two identity decompositions,
one on \(\spanop U_J\) and one on its orthogonal complement.  Repeating this
finite procedure yields an orthogonal direct-sum decomposition for the vectors
\(\{u_j:c_j>0\}\) into strict blocks.  Applying \(R^{-1}\) transfers it to a
direct-sum decomposition of the original spans \(V_{J_\ell}\).  Therefore
\(A\) is totally reducible for \(c\).

\medskip
\noindent\textbf{Step 3: total reducibility gives a solution.}
Conversely, suppose \(A\) is totally reducible for \(c\).  Discard the indices
with \(c_j=0\), and write the positive index set as a disjoint union
\[
  J_1\cup\cdots\cup J_r
\]
with
\[
  \R^n=V_{J_1}\oplus\cdots\oplus V_{J_r},
  \qquad
  c_{J_\ell}\in K_{A_{J_\ell}}^\circ.
\]
Choose an invertible linear map \(T\) such that the subspaces
\[
  E_\ell:=T V_{J_\ell}
\]
are mutually orthogonal.  The rank inequalities and their strictness are
preserved by \(T\), so
\[
  c_{J_\ell}\in K_{T A_{J_\ell}}^\circ
\]
inside \(E_\ell\).  By Step 1, for each \(\ell\) there is a positive definite
operator \(R_\ell:E_\ell\to E_\ell\) such that
\[
  \sum_{j\in J_\ell}c_j
  \left(\frac{R_\ell T a_j}{|R_\ell T a_j|}\right)\otimes
  \left(\frac{R_\ell T a_j}{|R_\ell T a_j|}\right)
  =\Id_{E_\ell}.
\]
Let \(R_0\) be the block-diagonal operator on the orthogonal sum
\(\R^n=E_1\oplus\cdots\oplus E_r\) whose restriction to \(E_\ell\) is
\(R_\ell\).  Then the invertible map \(L=R_0T\) satisfies
\[
  \sum_j c_j
  \left(\frac{La_j}{|La_j|}\right)\otimes
  \left(\frac{La_j}{|La_j|}\right)
  =\Id.
\]
Finally take the polar decomposition \(L=QS\), where \(Q\) is orthogonal and
\(S\) is symmetric positive definite.  Since \(La_j/|La_j|=Q(Sa_j/|Sa_j|)\),
conjugating the last identity by \(Q^*\) gives
\[
  \sum_j c_j
  \left(\frac{Sa_j}{|Sa_j|}\right)\otimes
  \left(\frac{Sa_j}{|Sa_j|}\right)
  =\Id.
\]
Thus the required invertible symmetric matrix exists.  This completes the
direct proof.
\end{proof}

\section*{Verification and Limitations}

The proof is purely finite-dimensional linear algebra plus the standard vector
matroid basis-polytope theorem.  It does not use entropy, Fisher information,
Brascamp--Lieb extremizers, or Gaussian optimization.  The theorem being proved
is not new: it is Theorem 4.3 of the source paper.  The packet answers the
authors' request for a direct proof in a purely linear-algebra context.

Bounded novelty check: local run indexes had no `0710.0870` packet.  Web search
for radial isotropic position, Tyler's M-estimator, Forster transforms, and the
normalized identity equation found close known formulations, especially the
subspace-condition criterion for Tyler's M-estimator and radial isotropic
position.  Therefore the novelty confidence is only for this compact direct
derivation and its explicit connection to the source theorem, not for the
criterion itself.

\section*{Human Review Recommendation}

Review as a likely valid full answer to the source's direct-proof request.
The main items to check are: the identification of \(K_A\) with the vector
matroid basis polytope, the existence of a minimizer of \(F\) in the relative
interior case, and the transfer from an arbitrary invertible map to a symmetric
positive map via polar decomposition.

\begin{thebibliography}{9}

\bibitem{CCE}
E. A. Carlen and D. Cordero-Erausquin,
\emph{Subadditivity of the Entropy and its Relation to Brascamp--Lieb Type
Inequalities}, arXiv:0710.0870, 2007.

\bibitem{Edmonds}
J. Edmonds,
\emph{Submodular functions, matroids, and certain polyhedra},
in \emph{Combinatorial Structures and their Applications}, Gordon and Breach,
1970, pp. 69--87.

\end{thebibliography}

\end{document}
